\section{Related Work}

Our work sits at the intersection of three research threads: controlled pre-training infrastructure, domain mixing optimization, and task-specific data selection. We review each and explain why, despite rich progress in each area, the per-domain $\times$ per-benchmark specificity question remains open.

\subsection{Controlled Pre-Training Infrastructure}

The Pile \cite{Gao2020Pile} introduced a 22-domain heterogeneous corpus designed to improve cross-domain generalization. With 800GB of text spanning Wikipedia, Books, GitHub, PubMed, StackExchange, and 17 other domains, The Pile remains the most carefully documented large-scale pre-training corpus in terms of domain composition and proportions.

Pythia \cite{Biderman2023Pythia} extended this infrastructure by training 16 autoregressive language models (70M--12B parameters) on The Pile with exact dataloader documentation, releasing 154 intermediate checkpoints per model. Pythia's explicit design goal was to enable training dynamics and data attribution studies --- yet prior work using Pythia has focused on memorization, scaling laws, and deduplication effects \cite{Biderman2023Pythia}, not on exploiting the checkpoint trajectory for domain exposure analysis. MiniPile \cite{Kaddour2023MiniPile} demonstrated that a 6GB embedding-filtered subset of The Pile retains $\sim$98\% of GLUE performance --- confirming that quality filtering matters --- but focused on aggregate performance rather than per-domain $\times$ per-benchmark decomposition.

We use the same Pythia/Pile infrastructure but exploit it differently: rather than comparing different training runs, we treat the 154-checkpoint trajectory as a panel dataset with naturally occurring domain exposure variation, enabling within-family analysis without new training.

\subsection{Domain Mixing Optimization}

The most direct precursor to our work is the line of research optimizing domain mixing ratios during pre-training. DoReMi \cite{DoReMi2023} uses a small proxy model to find domain weights that minimize excess loss, demonstrating that optimized domain mixing improves aggregate validation loss and downstream benchmarks. RegMix \cite{Liu2024RegMix} extends this to regression-based mixing law: fitting a linear model from domain weights to validation loss at proxy-model scale, then extrapolating to larger models. This linear mixing law framework is the statistical foundation our panel regression adopts, though we apply it to individual benchmark scores rather than aggregate validation loss.

Data Mixing Laws \cite{Ye2024DataMixingLaws} provided the most direct evidence that domain mixing ratios have predictable effects: their scaling laws incorporate domain proportions as explicit predictors. Our work is consistent with this framing but adds within-family variation as the source of domain exposure differences, rather than cross-run variation. None of these approaches produce a per-domain coefficient for individual benchmarks (MMLU, HellaSwag, ARC, WinoGrande separately). They optimize or characterize \textit{aggregate} performance. Our panel regression framework is the first designed to estimate benchmark-specific domain coefficients from within-family variation.

\subsection{Task-Specific Data Selection}

CoLoR-Filter \cite{Brandfonbrener2024CoLoR} demonstrates that task-conditioned data selection achieves 11--25$\times$ data efficiency for specific downstream benchmarks --- the theoretical motivation for domain-benchmark specificity. Topic Over Source \cite{Peng2025TopicOverSource} provides the finding that topic-based mixing consistently outperforms source-based mixing, suggesting that the \textit{semantic content} of training data matters more than the source label --- consistent with our cognitive content proxy approach.

\subsection{Benchmark Contamination}

LatestEval \cite{Li2023LatestEval} and Evading Data Contamination Detection \cite{Dekoninck2024Evading} highlight the difficulty of contamination auditing. We include contamination-adjusted scores (lm-evaluation-harness 13-gram decontamination) and document this as a limitation.
